\vspace{0.5em}\noindent\textbf{{[}FCAJ2026{]} What is AWS Bedrock Knowledge Bases? Why is it
the perfect ``missing piece'' for Serverless RAG architecture?}

\begin{center}\rule{0.5\linewidth}{0.5pt}\end{center}

\vspace{0.5em}\noindent\textbf{Introduction}

After learning how to build Medical AI Assistant systems using the RAG
(Retrieval-Augmented Generation) approach from scratch, I began facing a
series of difficulties in building Data Pipelines on my own. What
surprised me was that when reading AWS official documentation and sample
solutions, most architectures utilized ``AWS Bedrock Knowledge Bases''
rather than providing instructions to deploy standalone Vector
Databases.

This made me wonder: ``If setting up custom chunking and a vector
database is the standard way to do RAG, why did AWS invent Bedrock
Knowledge Bases?''

After reading the official documentation and directly deploying this
service for a medical assistant project, I realized that Bedrock
Knowledge Bases is not just a ``vector storage place'' but a fully
managed service that oversees the entire RAG workflow, turning every
complex stage into Serverless.

\begin{center}\rule{0.5\linewidth}{0.5pt}\end{center}

\vspace{0.5em}\noindent\textbf{Self-managing RAG Infrastructure - An Operational
Nightmare}

Previously, I thought building a RAG system mainly revolved around
writing scripts. For instance, for the system to read tens of thousands
of medical PDF or CSV pages, I had to manually write code to chunk the
text, call embedding model APIs to convert the text, and then self-host
a Vector Database like Milvus or Qdrant.

This approach is quite fun for small projects. However, as the volume of
medical data increases, managing this infrastructure becomes difficult.
Maintaining a stable Vector Database or writing cronjobs/Airflow
pipelines to continuously synchronize data generates numerous
operational (Ops) issues.

This is why modern cloud systems encourage the use of Managed Serverless
services - a method of offloading the entire burden of managing
underlying infrastructure to the cloud provider.

\begin{center}\rule{0.5\linewidth}{0.5pt}\end{center}

\vspace{0.5em}\noindent\textbf{What is AWS Bedrock Knowledge Bases?}

According to official AWS documentation, ``Knowledge bases for Amazon
Bedrock'' is a fully managed capability that helps connect Foundation
Models with enterprise internal data sources to perform RAG.

In other words, AWS Bedrock Knowledge Bases fully automates the data
processing pipeline. Instead of having to do every step manually, this
service handles:

\begin{enumerate}
\def\labelenumi{\arabic{enumi}.}
\tightlist
\item
  \textbf{Automated Chunking}: Semantically splitting complex medical
  documents.
\item
  \textbf{Automated Embedding}: Transforming text into vectors using
  models like Amazon Titan or Cohere.
\item
  \textbf{Automated Indexing}: Storing and indexing data in a hidden
  Serverless Vector Database backend (such as Amazon OpenSearch
  Serverless).
\end{enumerate}

This is the point that changed my perspective on building AI
applications. Previously, I thought AI engineers had to manually control
every step of the Data Pipeline. However, after using Bedrock Knowledge
Bases, I realized that ``infrastructure operational work should be
automated to focus on application logic''.

\begin{center}\rule{0.5\linewidth}{0.5pt}\end{center}

\vspace{0.5em}\noindent\textbf{Workflow with AWS Bedrock Knowledge Bases}

During the system building process, integration and operation became
incredibly simple:

\textbf{Upload Data to Amazon S3} \emph{Function}: Amazon S3 acts as the
raw Data Source. Clinical documents and treatment protocols simply need
to be dropped into an S3 Bucket.

\textbf{Start Ingestion Job (Data Synchronization)} \emph{Function}:
With just a click on the AWS Console or an API call, the system
automatically synchronizes new documents from S3, handles chunking and
embedding to update the AI's knowledge base without causing any
downtime.

\textbf{Integration via LangChain} \emph{Function}: Allows Python
applications to retrieve context easily. Instead of writing complex
database connection code, I only need to use
\texttt{AmazonKnowledgeBasesRetriever} combined with the Knowledge Base
ID to query (including Hybrid Search capabilities).

\begin{center}\rule{0.5\linewidth}{0.5pt}\end{center}

\vspace{0.5em}\noindent\textbf{The Role of Amazon S3 in the Architecture}

Initially, I only viewed Amazon S3 as a place to store static files.
However, when combined with Bedrock Knowledge Bases, I realized S3 is
the ``knowledge gateway'' of the entire system. Whenever the hospital
issues a new medication list or protocol, the administrator simply
pushes the file to S3. With no code intervention and no application
redeployment required, the AI system instantly gains the latest
knowledge.

It is safe to say that Amazon S3 combined with Bedrock Knowledge Bases
creates an automated, infinite stream of knowledge updates.

\begin{center}\rule{0.5\linewidth}{0.5pt}\end{center}

\vspace{0.5em}\noindent\textbf{Why is AWS Bedrock Knowledge Bases the ``true love'' for
AI Engineers?}

This is the conclusion I reached after completing the project. The
biggest difference lies in resource liberation.

If building RAG manually, an engineer has to play the role of an entire
DevOps team: struggling with infrastructure, fixing pipeline errors, and
monitoring the Vector Database. Conversely, Bedrock Knowledge Bases
shifts all that burden to AWS. This service handles all complex implicit
tasks and returns to us the cleanest, simplest APIs.

Consequently, AWS Bedrock Knowledge Bases helps engineering teams save
weeks of work, allowing them to channel their complete focus into
optimizing clinical business logic and improving patient experiences.

\begin{center}\rule{0.5\linewidth}{0.5pt}\end{center}

\vspace{0.5em}\noindent\textbf{References}

\begin{enumerate}
\def\labelenumi{\arabic{enumi}.}
\tightlist
\item
  AWS. ``Knowledge bases for Amazon Bedrock''.
  (\url{https://docs.aws.amazon.com/bedrock/latest/userguide/knowledge-base.html})
\item
  LangChain. ``Amazon Knowledge Bases Integrations''.
  (\url{https://python.langchain.com/docs/integrations/retrievers/bedrock/})
\end{enumerate}
